
%%FILE blueprint/src/asymptotic_defs.tex
\chapter{Basic asymptotic properties of functions}

\begin{definition}
    \label{def:asymp_pos}
    \lean{AsympPos}
    \leanok
    $f$ is asymptotically positive if there exists $x_0$ such that $f(x) > 0$
    for all $x \ge x_0$.

\end{definition}

\begin{definition}
    \label{def:asymp_neg}
    \lean{AsympNeg}
    \leanok
    $f$ is asymptotically negative if there exists $x_0$ such that $f(x) < 0$
    for all $x \ge x_0$.

\end{definition}

\begin{definition}
    \label{def:asymp_nonpos}
    \lean{AsympNonpos}
    \leanok
    $f$ is asymptotically nonpositive if there exists $x_0$ such that $f(x) \ge 0$
    for all $x \ge x_0$.

\end{definition}

\begin{definition}
    \label{def:asymp_nonneg}
    \lean{AsympNonneg}
    \leanok
    $f$ is asymptotically negative if there exists $x_0$ such that $f(x) \le 0$
    for all $x \ge x_0$.

\end{definition}

\begin{definition} 
    \label{def:asymp_le} 
    \lean{AsympLE}
    \leanok
    $f$ is asymptotically less than $g$ if there exists $x_0$ such that
    $f(x) \le g(x)$ for all $x \ge x_0$.

\end{definition}

\begin{definition}
    \label{def:asymp_ge}
    \lean{AsympGE}
    \leanok
    $f$ is asymptotically greater than $g$ if there exists $x_0$ such that
    $f(x) \ge g(x)$ for all $x \ge x_0$.

\end{definition}

\section{Asymptotic positivity and negativity}

\begin{lemma}
    \label{lemma:asymp_neg_of_pos}
    \lean{asymp_neg_of_pos}
    \leanok
    \uses{def:asymp_pos, def:asymp_neg}
    If $f$ is asymptotically negative, then $-f$ is asymptotically positive.
\end{lemma}

\begin{proof}
    \leanok
    By definition of asymptotic positivity, there exists an $x\_0$ such that $f(x) > 0$
    for all $x > x\_0$. It follows that $-f(x) > 0$, which is what is needed.
\end{proof}

\begin{lemma}
    \label{lemma:asymp_pos_of_neg}
    \lean{asymp_pos_of_neg}
    \leanok
    \uses{def:asymp_pos, def:asymp_neg}
    If $f$ is asymptotically positive, then $-f$ is asymptotically negative.
\end{lemma}

\begin{proof}
    \leanok
    By definition of asymptotic negativity, there exists an $x\_0$ such that $f(x) < 0$
    for all $x > x\_0$. It follows that $-f(x) < 0$, which is what is needed.
\end{proof}

\section{Asymptotic inequality}

\subsection{Positivity and negativity}

\begin{lemma}
    \label{lemma:asymp_pos_of_pos_le}
    \lean{asymp_pos_of_pos_le}
    \leanok
    \uses{def:asymp_pos, def:asymp_le}
    Let $f$ be asymptotically positive and let it be asymptotically less than $g$.
    Then $g$ is asymptotically positive.
\end{lemma}

\begin{proof}
    \leanok
    Assume with no loss of generality that $f(x) > 0$ for all $x > x\_0$ and that
    $f(y) \le g(y)$ for all $y > y\_0$. Let $z\_0 = \max \{ x\_0, y\_0 \}$. By transitivity
    of the inequality relations, we have $g(z) > 0$ for all $z > z\_0$. 
\end{proof}

\begin{lemma}
    \label{lemma:asymp_neg_of_le_neg}
    \lean{asymp_neg_of_le_neg}
    \leanok
    \uses{def:asymp_neg, def:asymp_le}
    Let $f$ be asymptotically less than $g$ and let $g$ be asymptotically negative.
    Then $f$ is asymptotically negative.
\end{lemma}

\begin{proof}
    \leanok
    Assume with no loss of generality that $g(x) < 0$ for all $x > x\_0$ and that
    $f(y) \le g(y)$ for all $y > y\_0$. Let $z\_0 = \max \{ x\_0, y\_0 \}$. By transitivity
    of the inequality relations, we have $f(z) < 0$ for all $z > z\_0$. 
\end{proof}

\begin{lemma}
    \label{lemma:asymp_pos_of_ge_pos}
    \lean{asymp_pos_of_ge_pos}
    \leanok
    \uses{def:asymp_pos, def:asymp_ge}
    Let $f$ be asymptotically greater than $g$ and let $g$ be asymptotically negative.
    Then $g$ is asymptotically negative.
\end{lemma}

\begin{proof}
    \leanok
    \uses{thm:asymp_le_iff_asymp_ge, lemma:asymp_pos_of_pos_le}
    By Lemma \ref{thm:asymp_le_iff_asymp_ge}, this statement is equivalent to Lemma
    \ref{lemma:asymp_pos_of_pos_le}.
\end{proof}

\begin{lemma}
    \label{lemma:asymp_neg_of_neg_ge}
    \lean{asymp_neg_of_neg_ge}
    \leanok
    \uses{def:asymp_neg, def:asymp_ge}
    Let $f$ be asymptotically negative and let it be asymptotically greater than $g$.
    Then $g$ is asymptotically negative.
\end{lemma}

\begin{proof}
    \leanok
    \uses{thm:asymp_le_iff_asymp_ge, lemma:asymp_neg_of_le_neg}
    By Theorem \ref{thm:asymp_le_iff_asymp_ge}, this statement is equivalent to Lemma
    \ref{lemma:asymp_neg_of_le_neg}.
\end{proof}


\subsection{Reflexivity}

\begin{lemma}
    \label{lemma:asymp_le_refl}
    \lean{asymp_le_refl}
    \leanok
    \uses{def:asymp_le}
    $f$ is asymptotically less than $f$.
\end{lemma}

\begin{proof}
    \leanok
    By reflexivity of $\le$, we have $f(x) \le f(x)$ for any given $x$.
\end{proof}

\begin{lemma}
    \label{lemma:asymp_ge_refl}
    \lean{asymp_ge_refl}
    \leanok
    \uses{def:asymp_ge}
    $f$ is asymptotically greater than $f$.
\end{lemma}

\begin{proof}
    \leanok
    By reflexivity of $\ge$, we have $f(x) \ge f(x)$ for any given $x$.
\end{proof}


\subsection{Equivalence}

\begin{lemma}
    \label{lemma:asymp_le_of_asymp_ge}
    \lean{asymp_le_of_asymp_ge}
    \leanok
    \uses{def:asymp_le, def:asymp_ge}
    Let $f$ be asymptotically less than $g$. Then $g$ is asymptotically
    greater than $f$.
\end{lemma}

\begin{proof}
    \leanok
    Since $f(x) \le g(x)$ for all $x > x\_0$, we have $g(x) \ge f(x)$.
\end{proof}

\begin{lemma}
    \label{lemma:asymp_ge_of_asymp_le}
    \lean{asymp_ge_of_asymp_le}
    \leanok
    \uses{def:asymp_le, def:asymp_ge}
    Let $f$ be asymptotically greater than $g$. Then $g$ is asymptotically
    less than $f$.
\end{lemma}

\begin{proof}
    \leanok
    Since $f(x) \ge g(x)$ for all $x > x\_0$, we have $g(x) \le f(x)$.
\end{proof}

\begin{theorem}
    \label{thm:asymp_le_iff_asymp_ge}
    \lean{asymp_le_iff_asymp_ge}
    \leanok
    \uses{def:asymp_le, def:asymp_ge}
    $f$ is asymptotically less than $g$ if and only if $g$ is asymptotically
    greater than $f$.
\end{theorem}

\begin{proof}
    \leanok
    \uses{lemma:asymp_le_of_asymp_ge, lemma:asymp_ge_of_asymp_le}
    Lemma \ref{lemma:asymp_le_of_asymp_ge} and Lemma \ref{lemma:asymp_ge_of_asymp_le} are both directions
    respectively.
\end{proof}


\subsection{Transitivity}

\begin{lemma}
    \label{lemma:asymp_le_trans}
    \lean{asymp_le_trans}
    \leanok
    If $f$ is asymptotically less than $g$ and $g$ is asymptotically less than $h$,
    then $f$ is asymptotically less than $h$.
\end{lemma}

\begin{proof}
    \leanok
    By assumption, $f(x) \le g(x)$ for all $x \ge x\_0$ and $g(y) \le h(y)$ for all
    $y \ge y\_0$. Let $z\_0 = \max \{ x\_0, y\_0 \}$. By transitivity, we have 
    $f(z) \le g(z)$ for all $z \ge z\_0$.
\end{proof}

\begin{lemma}
    \label{lemma:asymp_ge_trans}
    \lean{asymp_ge_trans}
    \leanok
    If $f$ is asymptotically greater than $g$ and $g$ is asymptotically greater than $h$,
    then $f$ is asymptotically greater than $h$.
\end{lemma}

\begin{proof}
    \leanok
    \uses{thm:asymp_le_iff_asymp_ge, lemma:asymp_le_trans}
    By the equivalence given by Lemma \ref{thm:asymp_le_iff_asymp_ge}, we can apply
    Lemma \ref{lemma:asymp_le_trans} in reverse, since $h$ is asymptotically less than
    $g$ and $g$ is asymptotically less than $f$.
\end{proof}


\subsection{Additivity}

\begin{lemma}
    \label{lemma:asymp_le_add}
    \lean{asymp_le_add}
    \leanok
    \uses{def:asymp_le}
    Let $f\_1$ be asymptotically less than $g\_1$ and $f\_2$ be asymptotically less
    than $g\_2$. Then $f\_1 + f\_2$ is asymptotically less than $g\_1 + g\_2$.
\end{lemma}

\begin{proof}
    \leanok
    Let $x\_0$ be such that $f\_1(x) \le g\_1(x)$ for all $x > x\_0$ and let $y\_0$ be
    such that $f\_2(y) \le g\_2(y)$ for all $y > y\_0$. Those exist due to assumptions.
    Now let $z\_0 = \max \{ x\_0, y\_0 \}$. By transitivity, $f\_1(z) \le g\_1(z)$ and 
    $f\_2(z) \le g\_2(z)$ for all $z > z\_0$. By additivity, we can merge both 
    inequalities by adding both terms on the left side and both terms on the right side.
    We thus get $f\_1(z) + f\_2(z) \le g\_1(z) + g\_2(z)$, which by definition and 
    and extensionality means that $f\_1 + f\_2$ is asymptotically less than $g\_1 + g\_2$.
\end{proof}

\begin{lemma}
    \label{lemma:asymp_ge_add}
    \lean{asymp_ge_add}
    \leanok
    \uses{def:asymp_ge}
    Let $f\_1$ be asymptotically greater than $g\_1$ and $f\_2$ be asymptotically greater
    than $g\_2$. Then $f\_1 + f\_2$ is asymptotically greater than $g\_1 + g\_2$.
\end{lemma}

\begin{proof}
    \leanok
    \uses{thm:asymp_le_iff_asymp_ge, def:asymp_le, lemma:asymp_le_add}
    By Theorem \ref{thm:asymp_le_iff_asymp_ge}, $g\_1$ and $g\_2$ are asymptotically less 
    than $f\_1$ and $f\_2$ respectively. It suffices to show that $g\_1 + g\_2$ is
    asymptotically less than $f\_1 + f\_2$, which is precisely the result of
    Lemma \ref{lemma:asymp_le_add}.
\end{proof}

\begin{lemma}
    \label{lemma:asymp_ge_add_pos}
    \lean{asymp_ge_add_pos}
    \leanok
    \uses{def:asymp_pos, def:asymp_ge}
    Let $f\_1$ be asymptotically positive. Let also $f\_2$ be asymptotically greater 
    than $g$. Then $f\_1 + f\_2$ is asymptotically greater than $g$.
\end{lemma}

\begin{proof}
    \leanok
    By definition, there exists some $x\_0$ such that $f\_1(x) > 0$ for all $x > x\_0$.
    We also have $f\_2(y) \ge g(y)$ for all $y > y\_0$ for some $y\_0$. Let
    $z\_0 = \max \{ x\_0, y\_0, \}$. We now have, for all $z > z\_0$ both $f\_1(z) > 0$
    and $f\_2(z) \ge g(z)$. By additivity, we have $f\_1(z) + f\_2(z) \ge g(z)$. 
\end{proof}

\begin{lemma}
    \label{lemma:asymp_le_add_neg}
    \lean{asymp_le_add_neg}
    \leanok
    \uses{def:asymp_neg, def:asymp_ge}
    Let $f\_1$ be asymptotically negative. Let also $f\_2$ be asymptotically less 
    than $g$. Then $f\_1 + f\_2$ is asymptotically less than $g$.
\end{lemma}

\begin{proof}
    \leanok
    By definition, there exists some $x\_0$ such that $f\_1(x) < 0$ for all $x > x\_0$.
    We also have $f\_2(y) \le g(y)$ for all $y > y\_0$ for some $y\_0$. Let
    $z\_0 = \max \{ x\_0, y\_0, \}$. We now have, for all $z > z\_0$ both $f\_1(z) < 0$
    and $f\_2(z) \le g(z)$. By additivity, we have $f\_1(z) + f\_2(z) \le g(z)$. 
\end{proof}


\subsection{Scalar multiplicativity}

\begin{lemma}
    \label{lemma:asymp_le_pos_smul}
    \lean{asymp_le_pos_smul}
    \leanok
    \uses{def:asymp_le}
    Let $c > 0$ and let $f$ be asymptotically less than $g$. Then $c \cdot f$ is asymptotically
    less than $c \cdot g$.
\end{lemma}

\begin{proof}
    \leanok
    This is a simple consequence of scalar multiplication by a positive constant.
\end{proof}

\begin{lemma}
    \label{lemma:asymp_ge_pos_smul}
    \lean{asymp_ge_pos_smul}
    \leanok
    \uses{def:asymp_ge}
    Let $c > 0$ and let $f$ be asymptotically greater than $g$. Then $c \cdot f$ is asymptotically
    greater than $c \cdot g$.
\end{lemma}

\begin{proof}
    \leanok
    \uses{thm:asymp_le_iff_asymp_ge, lemma:asymp_le_pos_smul}
    By applying Theorem \ref{thm:asymp_le_iff_asymp_ge}, the proof boils down to proving that
    $c \cdot g$ is asymptotically less than $c \cdot f$, which is precisely shown
    by Lemma \ref{lemma:asymp_le_pos_smul}.
\end{proof}

\begin{lemma}
    \label{lemma:asymp_le_neg_smul}
    \lean{asymp_le_neg_smul}
    \leanok
    \uses{def:asymp_neg, def:asymp_le, def:asymp_ge}
    Let $c < 0$ and let $f$ be asymptotically less than $g$. Then $c \cdot f$ is asymptotically
    greater than $c \cdot g$.
\end{lemma}

\begin{proof}
    \leanok
    This is a simple consequence of the fact that if $f(x) \le g(x)$, then for a $c < 0$
    we have $c \cdot f(x) \ge c \cdot g(x)$.
\end{proof}

\begin{lemma}
    \label{lemma:asymp_ge_neg_smul}
    \lean{asymp_ge_neg_smul}
    \leanok
    \uses{def:asymp_neg, def:asymp_ge, def:asymp_le}
    Let $c < 0$ and let $f$ be asymptotically greater than $g$. Then $c \cdot f$ is asymptotically
    less than $c \cdot g$.
\end{lemma}

\begin{proof}
    \leanok
    \uses{thm:asymp_le_iff_asymp_ge, lemma:asymp_le_neg_smul}
    Similar to above, the proof is a direct application of Theorem \ref{thm:asymp_le_iff_asymp_ge}
    and Lemma \ref{lemma:asymp_le_neg_smul}.
\end{proof}


\subsection{Multiplicativity}

\begin{theorem}
    \label{thm:asymp_le_nonneg_mul}
    \lean{asymp_le_nonneg_mul}
    \leanok
    \uses{def:asymp_nonneg, def:asymp_le}
    Let $f\_1$ and $f\_2$ be asymptotically nonnegative functions. If f$\_1$ is asymptotically
    less than $g\_1$ and $f\_2$ is asymptotically less than $g\_2$, then $f\_1 * f\_2$
    is asymptotically less than $g\_1 * g\_2$.
\end{theorem}

\begin{proof}
    \leanok
    Asymptotic properties give constants $x\_i, 0 \le i \le 3$, above which each property holds.
    We take $x_M = \max_{0 \le i \le 3} x_i$ as the constant of the needed property. 
    For all $x > x_M$, all given asymptotic properties hold, so the wanted property holds
    by properties of the inequality relation.
\end{proof}

\begin{theorem}
    \label{thm:asymp_ge_nonpos_mul}
    \lean{asymp_ge_nonpos_mul}
    \leanok
    \uses{def:asymp_nonpos, def:asymp_ge, def:asymp_le}
    Let $f\_1$ and $f\_2$ be asymptotically nonpositive functions. If f$\_1$ is asymptotically
    greater than $g\_1$ and $f\_2$ is asymptotically greater than $g\_2$, then $f\_1 * f\_2$
    is asymptotically less than $g\_1 * g\_2$.
\end{theorem}

\begin{proof}
    \leanok
    Analogously to above, the proof comes from taking the maximum of asymptotic constants 
    as the asymptotic lower bound for nonpositivity. This time, however, the inequality
    flips due to nonpositive terms in $f\_1(n) * f\_2(n) \le g\_1(n) * g\_2(n)$, since 
    $n$ is larger than all of the asymptotic lower bounds.
\end{proof}


%%FILE blueprint/src/asymptotic_growth.tex
\chapter{Asymptotic growth}

\begin{definition}
    \label{def:asymp_bounded_above}
    \lean{AsympBoundedAbove}
    \leanok
    \uses{def:asymp_le}
    $f$ is asymptotically bounded above by $g$ if there exists a $k > 0$ 
    such that $f$ is asymptotically less than $k*g$.
\end{definition}

\begin{definition}
    \label{def:asymp_bounded_below}
    \lean{AsympBoundedBelow}
    \leanok
    \uses{def:asymp_ge}
    $f$ is asymptotically bounded below by $g$ if there exists $k > 0$ 
    such that $f$ is asymptotically greater than $k*g$.
\end{definition}

\begin{definition}
    \label{def:asymp_bounded}
    \lean{AsympBounded}
    \leanok
    \uses{def:asymp_bounded_above, def:asymp_bounded_below}
    $f$ is asymptotically bounded by $g$ if $f$ is asymptotically bounded
    above and below by $g$.
\end{definition}

\begin{definition}
    \label{def:asymp_right_dom}
    \lean{AsympRightDom}
    \leanok
    \uses{def:asymp_le}
    $f$ is asymptotically dominated by $g$ if for all $k > 0$ $f$ is asymptotically 
    less than $k*g$.

\end{definition}

\begin{definition}
    \label{def:asymp_left_dom}
    \lean{AsympLeftDom}
    \leanok
    \uses{def:asymp_ge}
    $f$ asymptotically dominates $g$ if for all $k > 0$ $(x)$ is asymptotically
    greater than $k*g$.

\end{definition}

\begin{lemma}
    \label{lemma:asymp_bounded_above_of_right_dom}
    \lean{asymp_bounded_above_of_right_dom}
    \leanok
    \uses{def:asymp_right_dom, def:asymp_bounded_above}
    If $f$ is dominated by $g$, then it's also bounded above by $g$.
\end{lemma}

\begin{proof}
    \leanok 
    The definitions of $f$ being dominated and bounded above by $g$ only differ
    in the quantifier before $k$ at the very start (universal for the hypothesis, existential
    for the goal), so it suffices to use any positive value for $k$. We can use $1$. 
    The desired result then follows directly.
\end{proof}

\begin{lemma}
    \label{lemma:asymp_bounded_below_of_left_dom}
    \lean{asymp_bounded_below_of_left_dom}
    \leanok
    \uses{def:asymp_left_dom, def:asymp_bounded_below}
    If $f$ dominates $g$, then it's bounded below by $g$.
\end{lemma}

\begin{proof}
    \leanok
    The proof is entirely analogous to the previous proof.
\end{proof}

\begin{lemma}
    \label{lemma:asymp_bounded_iff}
    \lean{asymp_bounded_iff}
    \leanok
    \uses{def:asymp_bounded_above, def:asymp_bounded_below, def:asymp_bounded}
    $f$ is asymptotically bounded above and below by $g$ if and only if $f$ is 
    asymptotically bounded by $g$.
\end{lemma}

\begin{proof}
    \leanok
    Both directions follow directly from the definition of asymptotic boundedness. 
\end{proof}

\begin{lemma}
    \label{lemma:not_asymp_pos_bounded_below_and_right_dom}
    \lean{not_asymp_pos_bounded_below_and_right_dom}
    \leanok
    \uses{def:asymp_pos, def:asymp_bounded_below, def:asymp_right_dom}
    Let $g$ be asymptotically positive. Then $f$ is not both asymptotically bounded 
    below by $g$ and asymptotically dominated by $g$.
\end{lemma}

\begin{proof}
    \leanok
    Suppose $f$ is asymptotically bounded below by $g$ and also asymptotically dominated
    by $g$. We need to find a contradiction to prove the statement. 

    First, we claim that there exists $x\_0$ such that for all $x \ge x\_0$ we have 
    $f(x) \ge k \cdot g(x)$ for some $k > 0$, $g(x) \ge 0$ and $f(x) \le k' \cdot g(x)$ 
    for all $k' > 0$. Each of the asymptotic assumption gives one such constant, so
    taking the maximum of all three gives the needed value.

    For the contradiction, consider $f(x) \le k' \cdot g(x)$ when $k' = k / 2$.
    In this case, we have $f(x) \le (k / 2) \cdot g(x)$, but we also have
    $k \cdot g(x) \le f(x)$, leading by transitivity to $k \cdot g(x) \le (k / 2) \cdot g(x)$,
    an obvious contradiction to the fact that $(k / 2) \cdot g(x) \le k \cdot g(x)$.
\end{proof}

\begin{theorem}
    \label{thm:not_asymp_pos_right_dom_of_bounded_below}
    \lean{not_asymp_pos_right_dom_of_bounded_below}
    \leanok
    \uses{def:asymp_pos, def:asymp_bounded_above, def:asymp_right_dom}
    Let $f$ be asymptotically positive. If $f$ is asymptotically bounded below by $g$, 
    then $f$ is not asymptotically dominated by $g$.
\end{theorem}

\begin{proof}
    \leanok
    \uses{lemma:not_asymp_pos_bounded_below_and_right_dom}
    This is a direct application of Lemma \ref{lemma:not_asymp_pos_bounded_below_and_right_dom}. 
\end{proof}

\begin{theorem}
    \label{thm:not_asymp_pos_bounded_below_of_right_dom}
    \lean{not_asymp_pos_bounded_below_of_right_dom}
    \leanok
    \uses{def:asymp_pos, def:asymp_bounded_above, def:asymp_right_dom}
    Let $f$ be asymptotically positive. If $f$ is asymptotically bounded below by $g$, 
    then $f$ is not asymptotically dominated by $g$.
\end{theorem}

\begin{proof}
    \leanok
    \uses{lemma:not_asymp_pos_bounded_below_and_right_dom}
    This is a direct application of Lemma \ref{lemma:not_asymp_pos_bounded_below_and_right_dom}. 
\end{proof}

\begin{lemma}
    \label{lemma:not_asymp_pos_bounded_above_and_left_dom}
    \lean{not_asymp_pos_bounded_above_and_left_dom}
    \leanok
    \uses{def:asymp_pos, def:asymp_bounded_above, def:asymp_left_dom}
    Let $g$ be asymptotically positive. Then it does not hold that both $f$ is asymptotically 
    bounded above by $g$ and $f$ asymptotically dominate $g$.
\end{lemma}

\begin{proof}
    \leanok
    The proof is analogous to the proof of Lemma 
    \ref{lemma:not_asymp_pos_bounded_below_and_right_dom}. This time, we set $k' = k + 1$
    and thus produce the false inequality $(k + 1) \cdot g(x) \le k \cdot g(x)$.
\end{proof}

\begin{theorem}
    \label{thm:not_asymp_left_dom_of_bounded_above_pos}
    \lean{not_asymp_left_dom_of_bounded_above_pos}
    \leanok
    \uses{def:asymp_pos, def:asymp_bounded_below, def:asymp_left_dom}
    Let $f$ be asymptotically positive. If $f$ is bounded below by $g$, then $f$ 
    does not dominate $g$.
\end{theorem}

\begin{proof}
    \leanok
    \uses{lemma:not_asymp_pos_bounded_above_and_left_dom}
    This is a direct application of Lemma \ref{lemma:not_asymp_pos_bounded_above_and_left_dom}. 
\end{proof}

\begin{theorem}
    \label{thm:not_asymp_bounded_above_of_left_dom_pos}
    \lean{not_asymp_bounded_above_of_left_dom_pos}
    \leanok
    \uses{def:asymp_pos, def:asymp_bounded_below, def:asymp_left_dom}
    Let $f$ be asymptotically positive. If $f$ is asymptotically bounded above by $g$, 
    then $f$ does not asymptotically dominate $g$.
\end{theorem}

\begin{proof}
    \leanok
    \uses{lemma:not_asymp_pos_bounded_above_and_left_dom}
    This is a direct application of Lemma \ref{lemma:not_asymp_pos_bounded_above_and_left_dom}. 
\end{proof}

\begin{lemma}
    \label{lemma:not_asymp_pos_left_and_right_dom}
    \lean{not_asymp_pos_left_and_right_dom}
    \leanok
    \uses{def:asymp_pos, def:asymp_left_dom, def:asymp_right_dom}
    Let $g$ be asymptotically positive. Then it is not true that both $f$ is asymptotically
    dominated by $g$ and that $f$ dominates $g$.
\end{lemma}

\begin{proof}
    Suppose $f$ both dominates $g$ and is dominated by $g$. Our goal is to find a
    contradiction. We have by definition the inequalities $g(x) > 0$, 
    $f(x) \ge k\_1 \cdot g(x)$ and $f(x) \le k\_2 \cdot g(x)$ for all $k\_1 > 0$, 
    $k\_2 > 0$ and for all $x \ge x\_0$ for some $x\_0$. Note that we use the same 
    constant $x\_0$ for all inequalities with no loss of generality. In fact, we shall also use the same $x\_0$ 
    in the asymptotic positivity condition $g(x) > 0$.

    Fix $k\_1 = 2$ and $k\_2 = 1$ (generally, we only need $k\_1 \ge k\_2$), so we now have
    $f(x) \ge 2 \cdot g(x)$ and $f(x) \le g(x)$. From these inequalities, it immediately
    follows that $2 \cdot g(x) \le g(x)$. However, since $1 \le 2$ and $g(x) > 0$, we have
    $g(x) < 2 \cdot g(x)$. We thus have two contradicting inequalities, finishing the proof.
\end{proof}

\begin{theorem}
    \label{thm:not_asymp_pos_left_dom_of_right_dom}
    \lean{not_asymp_pos_left_dom_of_right_dom}
    \leanok
    \uses{def:asymp_pos, def:asymp_left_dom, def:asymp_right_dom}
    Let $g$ be asymptotically positive. If $f$ is asymptotically dominated by $g$, then
    $f$ does not asymptotically dominate $g$.
\end{theorem}

\begin{proof}
    \leanok
    \uses{lemma:not_asymp_pos_left_and_right_dom}
    This is a direct application of Lemma \ref{lemma:not_asymp_pos_left_and_right_dom}. 
\end{proof}

\begin{theorem}
    \label{thm:not_asymp_pos_right_dom_of_left_dom}
    \lean{not_asymp_pos_right_dom_of_left_dom}
    \leanok
    \uses{def:asymp_pos, def:asymp_left_dom, def:asymp_right_dom}
    Let $g$ be asymptotically positive. If $f$ asymptotically dominates $g$, then
    $f$ is not asymptotically dominated by $g$.
\end{theorem}

\begin{proof}
    \leanok
    \uses{lemma:not_asymp_pos_left_and_right_dom}
    This is a direct application of Lemma \ref{lemma:not_asymp_pos_left_and_right_dom}. 
\end{proof}


\section{Reflexivity}

\begin{lemma}
    \label{lemma:asymp_bounded_refl}
    \lean{asymp_bounded_refl}
    \leanok
    \uses{def:asymp_bounded}
    $f$ is asymptotically bounded by itself. 

\end{lemma}

\begin{proof}
    \leanok
    \uses{def:asymp_bounded, def:asymp_bounded_above, def:asymp_bounded_below}
    Proving asymptotic boundedness is equivalent to proving boundedness above and below.
    Both can be proved the same way - we choose $1_K$ for $K$ and $1_R$ for $N$, then the
    required asymptotic growth properties follow from definitions of identity elements for 
    $K$ and $R$.
\end{proof}

\begin{lemma}
    \label{lemma:asymp_bounded_above_refl}
    \lean{asymp_bounded_above_refl}
    \leanok
    \uses{def:asymp_bounded_above}
    $f$ is asymptotically bounded above by itself.
\end{lemma}

\begin{proof}
    \leanok
    \uses{lemma:asymp_bounded_refl}
    This follows directly from \ref{lemma:asymp_bounded_refl}.
\end{proof}

\begin{lemma}
    \label{lemma:asymp_bounded_below_refl}
    \lean{asymp_bounded_below_refl}
    \leanok
    \uses{def:asymp_bounded_below}
    $f$ is asymptotically bounded below by itself.
\end{lemma}

\begin{proof}
    \leanok
    \uses{lemma:asymp_bounded_refl}
    This follows directly from \ref{lemma:asymp_bounded_refl}.
\end{proof}


\section{Transitivity}

\begin{lemma}
    \label{lemma:asymp_bounded_above_trans}
    \lean{asymp_bounded_above_trans}
    \leanok
    \uses{def:asymp_bounded_above}
    If $f$ is asymptotically bounded above by $g$ and $g$ is asymptotically bounded above
    by $h$, then $f$ is asymptotically bounded above by $h$.
\end{lemma}

\begin{proof}
    \leanok
    \uses{lemma:asymp_le_pos_smul, lemma:asymp_le_trans}
    Let $k\_1$ and $k\_2$ be the constants such that $f(x) \le k\_1 \cdot g(x)$ and
    $g(x) \le k\_2 \cdot h(x)$ for sufficiently large $x$. By multiplicativity and 
    transitivity, we have $f(x) \le k\_1 \cdot k\_2 \cdot h(x)$.
\end{proof}

\begin{lemma}
    \label{lemma:asymp_bounded_below_trans}
    \lean{asymp_bounded_below_trans}
    \leanok
    \uses{def:asymp_bounded_below}
    If $f$ is asymptotically bounded below by $g$ and $g$ is asymptotically bounded below
    by $h$, then $f$ is asymptotically bounded below by $h$.
\end{lemma}

\begin{proof}
    \leanok
    \uses{lemma:asymp_ge_pos_smul, lemma:asymp_ge_trans}
    Let $k\_1$ and $k\_2$ be the constants such that $f(x) \ge k\_1 \cdot g(x)$ and
    $g(x) \ge k\_2 \cdot h(x)$ for sufficiently large $x$. By multiplicativity and 
    transitivity, we have $f(x) \ge k\_1 \cdot k\_2 \cdot h(x)$.
\end{proof}

\begin{lemma}
    \label{lemma:asymp_bounded_trans}
    \lean{asymp_bounded_trans}
    \leanok
    \uses{def:asymp_bounded}
    If $f$ is asymptotically bounded by $g$ and $g$ is asymptotically bounded 
    by $h$, then $f$ is asymptotically bounded by $h$.
\end{lemma}

\begin{proof}
    \leanok
    \uses{lemma:asymp_bounded_above_trans, lemma:asymp_bounded_below_trans}
    A direct consequence of Lemma \ref{lemma:asymp_bounded_above_trans} and Lemma
    \ref{lemma:asymp_bounded_below_trans}.
\end{proof}

\begin{lemma}
    \label{lemma:asymp_right_dom_trans}
    \lean{asymp_right_dom_trans}
    \leanok
    \uses{def:asymp_right_dom}
    If $f$ is asymptotically dominated by $g$ and $g$ is asymptotically dominated by $h$, then 
    $f$ is asymptotically dominated by $h$.
\end{lemma}

\begin{proof}
    \leanok
    \uses{lemma:asymp_le_pos_smul, lemma:asymp_le_trans}
    Let $k > 0$. Then by first assumption, we have $f(x) \le k \cdot g(x)$ for large 
    $x$. By the second assumption, we have $g(x) \le h(x)$ for large $x$. By multiplicativity 
    and transitivity, we get $f(x) \le k \cdot h(x)$.
\end{proof}

\begin{lemma}
    \label{lemma:asymp_left_dom_trans}
    \lean{asymp_left_dom_trans}
    \leanok
    \uses{def:asymp_left_dom}
    If $f$ asymptotically dominates $g$ and $g$ asymptotically dominates $h$, then 
    $f$ asymptotically dominates $h$.
\end{lemma}

\begin{proof}
    \leanok
    \uses{lemma:asymp_ge_pos_smul, lemma:asymp_ge_trans}
    Let $k > 0$. Then by first assumption, we have $f(x) \ge k \cdot g(x)$ for large 
    $x$. By the second assumption, we have $g(x) \ge h(x)$ for large $x$. By multiplicativity 
    and transitivity, we get $f(x) \ge k \cdot h(x)$.
\end{proof}


\section{Scalar multiplicativity}

\begin{lemma}
    \label{lemma:asymp_bounded_above_pos_smul}
    \lean{asymp_bounded_above_pos_smul}
    \leanok
    \uses{def:asymp_bounded_above}
    Let $c > 0$. If $f$ is bounded above by $g$, then $c \cdot f$ is also bounded
    above by $g$.
\end{lemma}

\begin{proof}
    \leanok
    \uses{lemma:asymp_le_pos_smul}
    Let $k$ be the constant such that $f$ is asymptotically less than $k \cdot g$.
    By positive scalar multiplicativity of asymptotic inequality, $c \cdot f$ is 
    asymptotically less than $c \cdot k \cdot g$. Since $c \cdot k > 0$, this implies 
    that $c \cdot f$ is bounded above by $g$.
\end{proof}

\begin{lemma}
    \label{lemma:asymp_bounded_below_pos_smul}
    \lean{asymp_bounded_below_pos_smul}
    \leanok
    \uses{def:asymp_bounded_below}
    Let $c > 0$. If $f$ is bounded below by $g$, then $c \cdot f$ is also bounded
    below by $g$.
\end{lemma}

\begin{proof}
    \leanok
    \uses{lemma:asymp_ge_pos_smul}
    Let $k$ be the constant such that $f$ is asymptotically greater than $k \cdot g$.
    By positive scalar multiplicativity of asymptotic inequality, $c \cdot f$ is 
    asymptotically greater than $c \cdot k \cdot g$. Since $c \cdot k > 0$, this implies
    that $c \cdot f$ is bounded below by $g$.
\end{proof}

\begin{lemma}
    \label{lemma:asymp_bounded_pos_smul}
    \lean{asymp_bounded_pos_smul}
    \leanok
    \uses{def:asymp_bounded}
    Let $c > 0$. If $f$ is bounded by $g$, then $c \cdot f$ is also bounded by $g$.
\end{lemma}

\begin{proof}
    \leanok
    \uses{lemma:asymp_bounded_above_pos_smul, lemma:asymp_bounded_below_pos_smul,}
    Above boundedness is exactly Lemma \ref{lemma:asymp_bounded_above_pos_smul} and
    below boundedness is exactly Lemma \ref{lemma:asymp_bounded_below_pos_smul}.
\end{proof}

\begin{lemma}
    \label{lemma:asymp_bounded_above_neg_smul}
    \lean{asymp_bounded_above_neg_smul}
    \leanok
    \uses{def:asymp_bounded_above}
    Let $c < 0$. If $f$ is bounded above by $g$, then $c \cdot f$ is bounded
    above by $-g$.
\end{lemma}

\begin{proof}
    \leanok
    \uses{lemma:asymp_le_neg_smul}
    Let $k$ be the constant such that $f$ is asymptotically less than $k \cdot g$.
    By positive scalar multiplicativity of asymptotic inequality, $-c \cdot f$ is 
    asymptotically less than $-c \cdot k \cdot g$. Since $-c \cdot k > 0$, this implies
    that $c \cdot f$ is bounded above by $-g$.
\end{proof}

\begin{lemma}
    \label{lemma:asymp_bounded_below_neg_smul}
    \lean{asymp_bounded_below_neg_smul}
    \leanok
    \uses{def:asymp_bounded_below}
    Let $c < 0$. If $f$ is bounded below by $g$, then $c \cdot f$ is bounded
    below by $-g$.
\end{lemma}

\begin{proof}
    \leanok
    \uses{lemma:asymp_le_neg_smul}
    Let $k$ be the constant such that $f$ is asymptotically greater than $k \cdot g$.
    By positive scalar multiplicativity of asymptotic inequality, $-c \cdot f$ is 
    asymptotically greater than $-c \cdot k \cdot g$. Since $-c \cdot k > 0$, this implies
    that $c \cdot f$ is bounded below by $-g$.
\end{proof}

\begin{lemma}
    \label{lemma:asymp_bounded_neg_smul}
    \lean{asymp_bounded_neg_smul}
    \leanok
    \uses{def:asymp_bounded}
    Let $c < 0$. If $f$ is bounded by $g$, then $c \cdot f$ is bounded by $-g$.
\end{lemma}

\begin{proof}
    \leanok
    \uses{lemma:asymp_bounded_above_neg_smul, lemma:asymp_bounded_below_neg_smul}
    Above boundedness is exactly Lemma \ref{lemma:asymp_bounded_above_neg_smul} and
    below boundedness is exactly Lemma \ref{lemma:asymp_bounded_below_neg_smul}.
\end{proof}


\section{Additivity}

\begin{lemma}
    \label{lemma:asymp_bounded_above_add}
    \lean{asymp_bounded_above_add}
    \leanok
    \uses{def:asymp_bounded_above}
    Let $f\_1$ and $f\_2$ be bounded above by $g$. Then $f\_1 + f\_2$ is also bounded 
    above by $g$.
\end{lemma}

\begin{proof}
    \leanok
    \uses{lemma:asymp_le_add}
    Let $k\_1$ and $k\_2$ be the constants such that $f\_1$ is asymptotically less than $k\_1 \cdot g$ 
    and $f\_2$ is asymptotically less than $k\_2 \cdot g$. By additivity of asymptotic
    inequality, $f\_1 + f\_2$ is asymptotically less than $(k\_1 * k\_2) \cdot g$.
    It directly follows that $f\_1 + f\_2$ is asymptotically bounded above by $g$.
\end{proof}

\begin{lemma}
    \label{lemma:asymp_bounded_below_add}
    \lean{asymp_bounded_below_add}
    \leanok
    \uses{def:asymp_bounded_below}
    Let $f\_1$ and $f\_2$ be bounded below by $g$. Then $f\_1 + f\_2$ is also bounded 
    below by $g$.
\end{lemma}

\begin{proof}
    \leanok
    \uses{lemma:asymp_ge_add}
    Let $k\_1$ and $k\_2$ be the constants such that $f\_1$ is asymptotically greater than $k\_1 \cdot g$ 
    and $f\_2$ is asymptotically greater than $k\_2 \cdot g$. By additivity of asymptotic
    inequality, $f\_1 + f\_2$ is asymptotically greater than $(k\_1 * k\_2) \cdot g$.
    It directly follows that $f\_1 + f\_2$ is asymptotically bounded below by $g$.
\end{proof}

\begin{lemma}
    \label{lemma:asymp_bounded_add}
    \lean{asymp_bounded_above_add}
    \leanok
    \uses{def:asymp_bounded}
    Let $f\_1$ and $f\_2$ be bounded by $g$. Then $f\_1 + f\_2$ is also bounded by $g$.
\end{lemma}

\begin{proof}
    \leanok
    \uses{lemma:asymp_bounded_above_add, lemma:asymp_bounded_below_add}
    This is proved directly by Lemma \ref{lemma:asymp_bounded_above_add} and Lemma
    \ref{lemma:asymp_bounded_below_add}.
\end{proof}

\begin{lemma}
    \label{lemma:asymp_bounded_below_add_pos}
    \lean{asymp_bounded_below_add_pos}
    \leanok
    \uses{def:asymp_bounded_below, def:asymp_pos}
    Let $f\_1$ be bounded below by $g$ and let $f\_2$ be asymptotically positive. 
    Then $f\_1 + f\_2$ is bounded below by $g$.
\end{lemma}

\begin{proof}
    \leanok
    \uses{lemma:asymp_ge_add_pos}
    This property immediately follows from Lemma \ref{lemma:asymp_ge_add_pos}.
\end{proof}

\begin{lemma}
    \label{lemma:asymp_bounded_above_add_neg}
    \lean{asymp_bounded_above_add_neg}
    \leanok
    \uses{def:asymp_bounded_above, def:asymp_neg}
    Let $f\_1$ be bounded above by $g$ and let $f\_2$ be asymptotically negative. 
    Then $f\_1 + f\_2$ is bounded above by $g$.
\end{lemma}

\begin{proof}
    \leanok
    \uses{lemma:asymp_le_add_neg}
    This property immediately follows from Lemma \ref{lemma:asymp_le_add_neg}.
\end{proof}

\begin{lemma}
    \label{lemma:asymp_bounded_add_pos_above}
    \lean{asymp_bounded_add_pos_above}
    \leanok
    \uses{def:asymp_bounded, def:asymp_bounded_above, def:asymp_pos}
    Let $f\_1$ be bounded by $g$. Let also $f\_2$ be asymptotically positive and
    bounded above by $g$. Then $f\_1 + f\_2$ is bounded by $g$.
\end{lemma}

\begin{proof}
    \leanok
    \uses{lemma:asymp_bounded_above_add, lemma:asymp_bounded_below_add_pos}
    This property immediately follows from Lemma \ref{lemma:asymp_bounded_above_add}
    and Lemma \ref{lemma:asymp_bounded_below_add_pos}.
\end{proof}

\begin{lemma}
    \label{lemma:asymp_bounded_add_neg_below}
    \lean{asymp_bounded_add_neg_below}
    \leanok
    \uses{def:asymp_bounded, def:asymp_bounded_below, def:asymp_neg}
    Let $f\_1$ be bounded by $g$. Let also $f\_2$ be asymptotically negative and
    bounded below by $g$. Then $f\_1 + f\_2$ is bounded by $g$.
\end{lemma}

\begin{proof}
    \leanok
    \uses{lemma:asymp_bounded_below_add, lemma:asymp_bounded_above_add_neg}
    This property immediately follows from Lemma \ref{lemma:asymp_bounded_below_add}
    and Lemma \ref{lemma:asymp_bounded_above_add_neg}.
\end{proof}

\begin{lemma}
    \label{lemma:asymp_bounded_add_pos_right_dom}
    \lean{asymp_bounded_add_pos_right_dom}
    \leanok
    \uses{def:asymp_bounded, def:asymp_right_dom, def:asymp_pos}
    Let $f\_1$ be bounded by $g$. Let also $f\_2$ be asymptotically positive and
    let $f\_2$ be dominated by $g$. Then $f\_1 + f\_2$ is bounded by $g$.
\end{lemma}

\begin{proof}
    \leanok
    \uses{lemma:asymp_bounded_add_pos_above, lemma:asymp_bounded_above_of_right_dom}
    We prove this with an application of Lemma \ref{lemma:asymp_bounded_add_pos_above}
    on Lemma \ref{lemma:asymp_bounded_above_of_right_dom}.
\end{proof}

\begin{lemma}
    \label{lemma:asymp_bounded_add_neg_left_dom}
    \lean{asymp_bounded_add_neg_left_dom}
    \leanok
    \uses{def:asymp_bounded, def:asymp_bounded_below, def:asymp_neg}
    Let $f\_1$ be bounded by $g$. Let also $f\_2$ be asymptotically negative and
    let $f\_2$ dominate $g$. Then $f\_1 + f\_2$ is bounded by $g$.
\end{lemma}

\begin{proof}
    \leanok
    \uses{lemma:asymp_bounded_add_neg_below, lemma:asymp_bounded_below_of_left_dom}
    We prove this with an application of Lemma \ref{lemma:asymp_bounded_add_neg_below}
    on Lemma \ref{lemma:asymp_bounded_below_of_left_dom}.
\end{proof}


\section{Multiplicativity}

\begin{lemma}
    \label{lemma:asymp_bounded_above_nonneg_mul}
    \lean{asymp_bounded_above_nonneg_mul}
    \leanok
    \uses{def:asymp_bounded_above, def:asymp_nonneg}
    Let $f\_1$ and $f\_2$ be asymptotically nonnegative functions such that $f\_1$ 
    is asymptotically bounded above by $g\_1$ and $f\_2$ is asymptotically bounded 
    above by $g\_2$. Then $f\_1 * f\_2$ is asymptotically bounded above by $g\_1 * g\_2$.
\end{lemma}

\begin{proof}
    \leanok
    \uses{thm:asymp_le_nonneg_mul}
    Let $k\_1$ and $k\_2$ be the constants such that $f\_1$ is asymptotically less than
    $g\_1$ and $f\_2$ is asymptotically less than $g\_2$. By multiplicativity of asymptotic 
    inequality, $f\_1 * f\_2$ is asymptotically less than $k\_1 \cdot g\_1 * k\_2 \cdot g\_2$,
    which is equivalent to asymptotic above boundedness of $f\_1 * f\_2$ by $g\_1 * g\_2$.
\end{proof}

\begin{lemma}
    \label{lemma:asymp_bounded_below_nonpos_mul}
    \lean{asymp_bounded_below_nonpos_mul}
    \leanok
    \uses{def:asymp_bounded_below, def:asymp_nonpos}
    Let $f\_1$ and $f\_2$ be asymptotically nonpositive functions such that $f\_1$ 
    is asymptotically bounded below by $g\_1$ and $f\_2$ is asymptotically bounded 
    below by $g\_2$. Then $f\_1 * f\_2$ is asymptotically bounded below by $g\_1 * g\_2$.
\end{lemma}

\begin{proof}
    \leanok
    \uses{thm:asymp_ge_nonpos_mul}
    Let $k\_1$ and $k\_2$ be the constants such that $f\_1$ is asymptotically greaterthan
    $g\_1$ and $f\_2$ is asymptotically greater than $g\_2$. By multiplicativity of asymptotic 
    inequality, $f\_1 * f\_2$ is asymptotically less than $k\_1 \cdot g\_1 * k\_2 \cdot g\_2$,
    which is equivalent to asymptotic below boundedness of $f\_1 * f\_2$ by $g\_1 * g\_2$.
\end{proof}


%%FILE blueprint/src/bachman_landau.tex
\chapter{Bachman-Landau notation}

Bachman-Landau family of notations is the name of a few closely related notations used in 
algorithm analysis. The most famous of them is the so-called big-O notation. While
most formulations are defined on functions from naturals or reals to reals, we define 
them more generally - requirements for the types of the domain and codomain vary between 
different properties. However, all properties hold for functions from a linearly ordered 
commutative ring to a linearly ordered field. In the rest of this page, we let $R$ be a linearly 
ordered commutative ring and $F$ be a linearly ordered field. We will also use symbols $f$, 
$f\_1$, $f\_2$, $g$, $g\_1$ and $g\_2$ for functions $R \to F$. Also, we let $M$ be 
a right $R$-module, although often only a (distributive) left multiplicative 
action on $R$ is required.


\section{Asymptotic sets}

\begin{definition}(Big O notation)
    \label{def:big_o}
    \lean{O}
    \leanok
    \uses{def:asymp_bounded_above}
    $f(x) \in O(g(x))$ if it is asymptotically bounded above by $g(x)$.
\end{definition}

\begin{definition}(Big Omega notation)
    \label{def:big_omega}
    \lean{Ω}
    \leanok
    \uses{def:asymp_bounded_below}
    $f(x) \in \Omega(g(x))$ if it is asymptotically bounded below by $g(x)$.
\end{definition}

\begin{definition}(Big Theta notation)
    \label{def:theta}
    \lean{Θ}
    \leanok
    \uses{def:asymp_bounded}
    $f(x) \in \Theta(g(x))$ if it is asymptotically bounded by $g(x)$. 
\end{definition}

\begin{definition}(Small O notation)
    \label{def:small_o}
    \lean{o}
    \leanok
    \uses{def:asymp_right_dom}
    $f(x) \in o(g(x))$ if it is asymptotically dominated by $g(x)$.
\end{definition}

\begin{definition}(Small Omega notation)
    \label{def:small_omega}
    \lean{ω}
    \leanok
    \uses{def:asymp_left_dom}
    $f(x) \in \omega(g(x))$ if it asymptotically dominates $g(x)$.
\end{definition}


\section{Relations between asymptotic sets} 

\begin{lemma}
    \label{lemma:big_o_of_small_o}
    \lean{O_of_o}
    \leanok
    \uses{def:small_o, def:big_o}
    If $f(x) \in o(g(x))$, then $f(x) \in O(f(x))$.
\end{lemma}

\begin{proof}
    \leanok
    \uses{lemma:asymp_bounded_above_of_right_dom}
    Since $o(g(x))$ and $O(f(x))$, we can simply use Lemma 
    \ref{lemma:asymp_bounded_above_of_right_dom}.
\end{proof}

\begin{theorem}
    \label{thm:big_omega_of_small_omega}
    \lean{Ω_of_ω}
    \leanok
    \uses{def:small_omega, def:big_omega}
    If $f(x) \in \omega(g(x))$, then $f(x) \in \Omega(g(x))$.
\end{theorem}

\begin{proof}
    \leanok
    \uses{lemma:asymp_bounded_below_of_left_dom}
    The proof is a simple application of Theorem 
    \ref{lemma:asymp_bounded_below_of_left_dom}.
\end{proof}

\begin{theorem}
    \label{thm:big_o_and_omega_iff_theta}
    \lean{O_Ω_iff_Θ}
    \leanok
    \uses{def:big_o, def:big_omega, def:theta}
    $f(x) \in O(g(x))$ and $f(x) \in \Omega(g(x))$ if and only if 
    $f(x) \in \Theta(g(x))$.
\end{theorem}

\begin{proof}
    \leanok
    \uses{lemma:asymp_bounded_iff}
    Similarly to previous proofs, the proof is a direct application of Lemma 
    \ref{lemma:asymp_bounded_iff}.
\end{proof}

\begin{lemma}
    \label{lemma:not_asymp_pos_theta_and_small_o}
    \lean{not_asymp_pos_Θ_and_o}
    \leanok
    Let $g$ be asymptotically positive. Then $f(x) \in \Theta(g(x))$ and $f(x) \in o(g(x))$
    are not both true.
\end{lemma}

\begin{proof}
    \leanok
    \uses{lemma:not_asymp_pos_bounded_below_and_right_dom}
    A direct application of Lemma \ref{lemma:not_asymp_pos_bounded_below_and_right_dom}.
\end{proof}

\begin{lemma}
    \label{lemma:not_asymp_pos_small_o_of_theta}
    \lean{not_asymp_pos_o_of_Θ}
    \leanok
    Let $g$ be asymptotically positive. If $f(x) \in \Theta(g(x))$ then $f(x) \notin o(g(x))$.
\end{lemma}

\begin{proof}
    \leanok
    \uses{lemma:not_asymp_pos_theta_and_small_o}
    A direct application of Lemma \ref{lemma:not_asymp_pos_theta_and_small_o}.
\end{proof}

\begin{lemma}
    \label{lemma:not_asymp_pos_theta_of_small_o}
    \lean{not_asymp_pos_o_of_Ω}
    \leanok
    Let $g$ be asymptotically positive. If $f(x) \in o(g(x))$ then $f(x) \notin \Theta(g(x))$.
\end{lemma}

\begin{proof}
    \leanok
    \uses{lemma:not_asymp_pos_theta_and_small_o}
    A direct application of Lemma \ref{lemma:not_asymp_pos_theta_and_small_o}.
\end{proof}

\begin{lemma}
    \label{lemma:not_asymp_pos_o_of_big_omega}
    \lean{not_asymp_pos_o_of_Ω}
    \leanok
    Let $g$ be asymptotically positive. If $f(x) \in \Omega(g(x))$ then $f(x) \notin o(g(x))$.
\end{lemma}

\begin{proof}
    \leanok
    \uses{thm:not_asymp_pos_right_dom_of_bounded_below}
    A direct application of Lemma \ref{thm:not_asymp_pos_right_dom_of_bounded_below}.
\end{proof}

\begin{lemma}
    \label{lemma:not_asymp_pos_big_omega_of_o}
    \lean{not_asymp_pos_Ω_of_o}
    \leanok
    Let $g$ be asymptotically positive. If $f(x) \in o(g(x))$ then $f(x) \notin \Omega(g(x))$.
\end{lemma}

\begin{proof}
    \leanok
    \uses{thm:not_asymp_pos_bounded_below_of_right_dom}
    A direct application of Lemma \ref{thm:not_asymp_pos_bounded_below_of_right_dom}.
\end{proof}

\begin{lemma}
    \label{lemma:not_asymp_pos_theta_and_small_omega}
    \lean{not_asymp_pos_Θ_and_ω}
    \leanok
    Let $g$ be asymptotically positive. Then $f(x) \in \Theta(g(x))$ and 
    $f(x) \in \omega(g(x))$ are not both true.
\end{lemma}

\begin{proof}
    \leanok
    \uses{lemma:not_asymp_pos_bounded_above_and_left_dom}
    A direct application of Lemma \ref{lemma:not_asymp_pos_bounded_above_and_left_dom}.
\end{proof}

\begin{lemma}
    \label{lemma:not_asymp_pos_small_omega_of_theta}
    \lean{not_asymp_pos_ω_of_Θ}
    \leanok
    Let $g$ be asymptotically positive. If $f(x) \in \Theta(g(x))$ then $f(x) \notin \omega(g(x))$.
\end{lemma}

\begin{proof}
    \leanok
    \uses{lemma:not_asymp_pos_theta_and_small_omega}
    A direct application of Lemma \ref{lemma:not_asymp_pos_theta_and_small_omega}.
\end{proof}

\begin{lemma}
    \label{lemma:not_asymp_pos_theta_of_small_omega}
    \lean{not_asymp_pos_ω_of_Θ}
    \leanok
    Let $g$ be asymptotically positive. If $f(x) \in \omega(g(x))$ then $f(x) \notin \Theta(g(x))$.
\end{lemma}

\begin{proof}
    \leanok
    \uses{lemma:not_asymp_pos_theta_and_small_omega}
    A direct application of Lemma \ref{lemma:not_asymp_pos_theta_and_small_omega}.
\end{proof}

\begin{lemma}
    \label{lemma:not_asymp_pos_o_and_small_omega}
    \lean{not_asymp_pos_o_and_ω}
    \leanok
    Let $g$ be asymptotically positive. Then $f(x) \in o(g(x))$ and 
    $f(x) \in \omega(g(x))$ are not both true.
\end{lemma}

\begin{proof}
    \leanok
    \uses{lemma:not_asymp_pos_left_and_right_dom}
    A direct application of Lemma \ref{lemma:not_asymp_pos_left_and_right_dom}.
\end{proof}

\begin{lemma}
    \label{lemma:not_asymp_pos_small_omega_of_o}
    \lean{not_asymp_pos_ω_of_o}
    \leanok
    Let $g$ be asymptotically positive. If $f(x) \in o(g(x))$ then $f(x) \notin \omega(g(x))$.
\end{lemma}

\begin{proof}
    \leanok
    \uses{lemma:not_asymp_pos_o_and_small_omega}
    A direct application of Lemma \ref{lemma:not_asymp_pos_o_and_small_omega}.
\end{proof}

\begin{lemma}
    \label{lemma:not_asymp_pos_o_of_small_omega}
    \lean{not_asymp_pos_o_of_ω}
    \leanok
    Let $g$ be asymptotically positive. If $f(x) \in o(g(x))$ then $f(x) \notin \omega(g(x))$.
\end{lemma}

\begin{proof}
    \leanok
    \uses{lemma:not_asymp_pos_o_and_small_omega}
    A direct application of Lemma \ref{lemma:not_asymp_pos_o_and_small_omega}.
\end{proof}


\section{Reflexivity}

\begin{lemma}
    \label{lemma:theta_refl}
    \lean{Θ_refl}
    \leanok
    \uses{def:theta}
    $f(x) \in \Theta(f(x))$.

\end{lemma}

\begin{proof}
    \leanok
    \uses{lemma:asymp_bounded_refl}
    Direct consequence of Lemma \ref{lemma:asymp_bounded_refl}.
\end{proof}

\begin{lemma}
    \label{lemma:big_o_refl}
    \lean{O_refl}
    \leanok
    \uses{def:big_o}
    $f(x) \in O(f(x))$.
\end{lemma}

\begin{proof}
    \leanok
    \uses{lemma:asymp_bounded_above_refl}
    This follows directly from Lemma \ref{lemma:asymp_bounded_above_refl}.
\end{proof}

\begin{lemma}
    \label{lemma:big_omega_refl}
    \lean{Ω_refl}
    \leanok
    \uses{def:big_omega}
    $f(x) \in \Omega(f(x))$.
\end{lemma}

\begin{proof}
    \leanok
    \uses{lemma:asymp_bounded_below_refl}
    This follows directly from Lemma \ref{lemma:asymp_bounded_below_refl}.
\end{proof}


\section{Transitivity}

\begin{lemma}
    \label{lemma:big_o_trans}
    \lean{O_trans}
    \leanok
    \uses{def:big_o}
    If $f(x) \in O(g(x))$ and $g(x) \in O(h(x))$, then $f(x) \in O(h(x))$.
\end{lemma}

\begin{proof}
    \leanok
    \uses{lemma:asymp_bounded_above_trans}
    This follows directly from Lemma \ref{lemma:asymp_bounded_above_trans}.
\end{proof}

\begin{lemma}
    \label{lemma:big_omega_trans}
    \lean{Ω_trans}
    \leanok
    \uses{def:big_omega}
    If $f(x) \in \Omega(g(x))$ and $g(x) \in \Omega(h(x))$, then $f(x) \in \Omega(h(x))$.
\end{lemma}

\begin{proof}
    \leanok
    \uses{lemma:asymp_bounded_below_trans}
    This follows directly from Lemma \ref{lemma:asymp_bounded_below_trans}.
\end{proof}

\begin{lemma}
    \label{lemma:theta_trans}
    \lean{Θ_trans}
    \leanok
    \uses{def:theta}
    If $f(x) \in \Theta(g(x))$ and $g(x) \in \Theta(h(x))$, then $f(x) \in \Theta(h(x))$.
\end{lemma}

\begin{proof}
    \leanok
    \uses{lemma:asymp_bounded_trans}
    This follows directly from Lemma \ref{lemma:asymp_bounded_trans}.
\end{proof}

\begin{lemma}
    \label{lemma:small_o_trans}
    \lean{o_trans}
    \leanok
    \uses{def:small_o}
    If $f(x) \in o(g(x))$ and $g(x) \in o(h(x))$, then $f(x) \in o(h(x))$. 
\end{lemma}

\begin{proof}
    \leanok
    \uses{lemma:asymp_right_dom_trans}
    This follows directly from Lemma \ref{lemma:asymp_right_dom_trans}.
\end{proof}

\begin{lemma}
    \label{lemma:small_omega_trans}
    \lean{ω_trans}
    \leanok
    \uses{def:small_omega}
    If $f(x) \in \omega(g(x))$ and $g(x) \in \omega(h(x))$, then $f(x) \in \omega(h(x))$.
\end{lemma}

\begin{proof}
    \leanok
    \uses{lemma:asymp_left_dom_trans}
    This follows directly from Lemma \ref{lemma:asymp_left_dom_trans}.
\end{proof}


\section{Scalar multiplicativity}

\begin{lemma}
    \label{lemma:big_o_pos_smul}
    \lean{O_pos_smul}
    \leanok
    \uses{def:big_o}
    Let $c > 0$. If $f(x) \in O(g(x))$, then $c \cdot f(x) \in O(g(x))$.
\end{lemma}

\begin{proof}
    \leanok
    \uses{lemma:asymp_bounded_above_pos_smul}
    This follows directly from Lemma \ref{lemma:asymp_bounded_above_pos_smul}.
\end{proof}

\begin{lemma}
    \label{lemma:big_omega_pos_smul}
    \lean{Ω_pos_smul}
    \leanok
    \uses{def:big_omega}
    Let $c > 0$. If $f(x) \in \Omega(g(x))$, then $c \cdot f(x) \in \Omega(g(x))$ is also bounded
    below by $g$.
\end{lemma}

\begin{proof}
    \leanok
    \uses{lemma:asymp_bounded_below_pos_smul}
    This follows directly from Lemma \ref{lemma:asymp_bounded_below_pos_smul}.
\end{proof}

\begin{lemma}
    \label{lemma:theta_pos_smul}
    \lean{Θ_pos_smul}
    \leanok
    \uses{def:theta}
    Let $c > 0$. If $f(x) \in \Theta(g(x))$, then $c \cdot f(x) \in \Theta(g(x))$.
\end{lemma}

\begin{proof}
    \leanok
    \uses{lemma:asymp_bounded_pos_smul}
    This follows directly from Lemma \ref{lemma:asymp_bounded_pos_smul}.
\end{proof}

\begin{lemma}
    \label{lemma:big_o_neg_smul}
    \lean{O_neg_smul}
    \leanok
    \uses{def:big_o}
    Let $c < 0$. If $f(x) \in O(g(x))$, then $c \cdot f(x) \in O(-g(x))$.
\end{lemma}

\begin{proof}
    \leanok
    \uses{lemma:asymp_bounded_above_neg_smul}
    This follows directly from Lemma \ref{lemma:asymp_bounded_above_neg_smul}.
\end{proof}

\begin{lemma}
    \label{lemma:big_omega_neg_smul}
    \lean{Ω_neg_smul}
    \leanok
    \uses{def:big_omega}
    Let $c < 0$. If $f \in \Omega(g(x))$, then $c \cdot f(x) \in \Omega(-g(x))$.
\end{lemma}

\begin{proof}
    \leanok
    \uses{lemma:asymp_bounded_below_neg_smul}
    This follows directly from Lemma \ref{lemma:asymp_bounded_below_neg_smul}.
\end{proof}

\begin{lemma}
    \label{lemma:theta_neg_smul}
    \lean{Θ_neg_smul}
    \leanok
    \uses{def:theta}
    Let $c < 0$. If $f(x) \in \Theta(g(x))$, then $c \cdot f \in \Theta(-g(x))$.
\end{lemma}

\begin{proof}
    \leanok
    \uses{lemma:asymp_bounded_neg_smul}
    This follows directly from Lemma \ref{lemma:asymp_bounded_neg_smul}.
\end{proof}


\section{Additivity}

\begin{lemma}
    \label{lemma:big_o_add}
    \lean{O_add}
    \leanok
    \uses{def:big_o}
    Let $f\_1(x), f\_2(x) \in O(g(x))$. Then $f\_1(x) + f\_2(x) \in O(g(x))$.
\end{lemma}

\begin{proof}
    \leanok
    \uses{lemma:asymp_bounded_above_add}
    This follows directly from Lemma \ref{lemma:asymp_bounded_above_add}.
\end{proof}

\begin{lemma}
    \label{lemma:big_omega_add}
    \lean{Ω_add}
    \leanok
    \uses{def:big_omega}
    Let $f\_1(x), f\_2(x) \in \Omega(g(x))$. Then $f\_1(x) + f\_2(x) \in \Omega(g(x))$.
\end{lemma}

\begin{proof}
    \leanok
    \uses{lemma:asymp_bounded_below_add}
    This follows directly from Lemma \ref{lemma:asymp_bounded_below_add}.
\end{proof}

\begin{lemma}
    \label{lemma:theta_add}
    \lean{Θ_add}
    \leanok
    \uses{def:theta}
    Let $f\_1(x), f\_2(x) \in \Theta(g(x))$. Then $f\_1(x) + f\_2(x) \in \Theta(g(x))$.
\end{lemma}

\begin{proof}
    \leanok
    \uses{lemma:asymp_bounded_add}
    This follows directly from Lemma \ref{lemma:asymp_bounded_add}.
\end{proof}

\begin{lemma}
    \label{lemma:big_omega_add_pos}
    \lean{Ω_add_pos}
    \leanok
    \uses{def:big_omega, def:asymp_pos}
    Let $f\_1(x) \in \Omega(g(x))$ and let $f\_2$ be asymptotically positive. 
    Then $f\_1(x) + f\_2(x) \in \Omega(g(x))$.
\end{lemma}

\begin{proof}
    \leanok
    \uses{lemma:asymp_bounded_below_add_pos}
    This follows directly from Lemma \ref{lemma:asymp_bounded_below_add_pos}.
\end{proof}

\begin{lemma}
    \label{lemma:big_o_add_neg}
    \lean{O_add_neg}
    \leanok
    \uses{def:big_o, def:asymp_neg}
    Let $f\_1(x) \in O(g(x))$ and let $f\_2$ be asymptotically negative. 
    Then $f\_1(x) + f\_2(x) \in O(g(x))$.
\end{lemma}

\begin{proof}
    \leanok
    \uses{lemma:asymp_bounded_above_add_neg}
    This follows directly from Lemma \ref{lemma:asymp_bounded_above_add_neg}.
\end{proof}

\begin{lemma}
    \label{lemma:theta_add_pos_O}
    \lean{Θ_add_pos_O}
    \leanok
    \uses{def:theta, def:big_o, def:asymp_pos}
    Let $f\_1(x) \in \Theta(g(x))$. Let also $f\_2$ be asymptotically positive and
    $f\_2(x) \in O(g(x))$. Then $f\_1(x) + f\_2(x) \in \Theta(g(x))$.
\end{lemma}

\begin{proof}
    \leanok
    \uses{lemma:asymp_bounded_add_pos_above}
    This follows directly from Lemma \ref{lemma:asymp_bounded_add_pos_above}.
\end{proof}

\begin{lemma}
    \label{lemma:theta_add_neg_Omega}
    \lean{Θ_add_neg_Ω}
    \leanok
    \uses{def:theta, def:big_omega, def:asymp_neg}
    Let $f\_1(x) \in \Theta(g(x))$. Let also $f\_2$ be asymptotically negative
    and $f\_2(x) \in \Omega(g(x))$. Then $f\_1(x) + f\_2(x) \in \Theta(g(x))$.
\end{lemma}

\begin{proof}
    \leanok
    \uses{lemma:asymp_bounded_add_neg_below}
    This follows directly from Lemma \ref{lemma:asymp_bounded_add_neg_below}.
\end{proof}

\begin{lemma}
    \label{lemma:theta_add_pos_o}
    \lean{Θ_add_pos_o}
    \leanok
    \uses{def:theta, def:small_o, def:asymp_pos}
    Let $f\_1(x) \in \Theta(g(x))$. Let also $f\_2$ be asymptotically positive and
    $f\_2(x) \in o(g(x))$. Then $f\_1(x) + f\_2(x) \in \Theta(g(x))$.
\end{lemma}

\begin{proof}
    \leanok
    \uses{lemma:asymp_bounded_add_pos_right_dom}
    This follows directly from Lemma \ref{lemma:asymp_bounded_add_pos_right_dom}.
\end{proof}

\begin{theorem}
    \label{thm:theta_add_neg_omega}
    \lean{Θ_add_neg_ω}
    \leanok
    \uses{def:theta, def:small_omega, def:asymp_neg}
    Let $f\_1(x) \in \Theta(g(x))$. Let also $f\_2$ be asymptotically negative and
    $f\_2(x) \in \omega(g(x))$. Then $f\_1(X) + f\_2(x) \in \Theta(g(x))$.
\end{theorem}

\begin{proof}
    \leanok
    \uses{lemma:asymp_bounded_add_neg_left_dom}
    This follows directly from Lemma \ref{lemma:asymp_bounded_add_neg_left_dom}.
\end{proof}


\section{Multiplicativity}

\begin{lemma}
    \label{lemma:big_o_nonneg_mul}
    \lean{O_nonneg_mul}
    \leanok
    \uses{def:big_o, def:asymp_nonneg}
    Let $f\_1$ and $f\_2$ be asymptotically nonnegative functions such that 
    $f\_1(x) \in O(g\_1(x))$ and $f\_2(x) \in O(g\_2(x))$. Then 
    $f\_1(x) * f\_2(x) \in O(g\_1(x) * g\_2(x))$.
\end{lemma}

\begin{proof}
    \leanok
    \uses{lemma:asymp_bounded_above_nonneg_mul}
    This follows directly from Lemma \ref{lemma:asymp_bounded_above_nonneg_mul}.
\end{proof}

\begin{lemma}
    \label{lemma:big_omega_nonpos_mul}
    \lean{Ω_nonpos_mul}
    \leanok
    \uses{def:big_omega, def:asymp_nonpos}
    Let $f\_1$ and $f\_2$ be asymptotically nonpositive functions such that 
    $f\_1(x) \in \Omega(g\_1(x))$ and $f\_2(x) \in \Omega(g\_2(x)$. 
    Then $f\_1(x) * f\_2(x) \in \Omega(g\_1(x) * g\_2(x))$.
\end{lemma}

\begin{proof}
    \leanok
    \uses{lemma:asymp_bounded_below_nonpos_mul}
    This follows directly from Lemma \ref{lemma:asymp_bounded_below_nonpos_mul}.
\end{proof}


%%FILE blueprint/src/bibliography.tex
\chapter{Bibliography}

\bibliographystyle{plain}
\bibliography{bibliography}


%%FILE blueprint/src/content.tex
% If you want to split the blueprint content into several files then
% the current file can be a simple sequence of \input. Otherwise It
% can start with a \section or \chapter for instance.

\title{Bachman-Landau notation and the Master theorem for 
        divide-and-conquer recurrences}

This project's goal is to formalize a proof of the Master theorem for 
divide-and-conquer recurrences. Along with the Master theorem, we also
state definitions and some properties of the Big-O and its sibling notations.

We use original and standard definitions of the Big-O notation 
\cite{bachmann1894, landau1909, sipser1996}, as well as Knuth's definition of 
Big-Omega and Big-Theta \cite{knuth1976}. We use a slightly adapted formulation 
of the Master Theorem from Cormen et al. \cite{cormen2022}, with the difference 
being due to our use of functions on natural numbers as opposed to their
more general formulation. The idea of the formalized proof is also adapted
from the same source, but their proof skips many technical steps, so our
proof appears to hold little resemblance.

\input{asymptotic_defs.tex}

\input{asymptotic_growth.tex}

\input{bachman_landau.tex}

\input{geom_sums.tex}

\input{master_theorem.tex}

\input{bibliography.tex}


%%FILE blueprint/src/geom_sums.tex
\chapter{Geometric sums}

Before we finally state and prove the Master theorem, we need to prove two
basic properties regarding geometric sums. These are both straightforward 
proofs with no clever tricks or surprises, but are still stated for the sake 
of completeness.

\begin{definition}
    \label{def:geom_sum}
    \lean{GeometricSum}
    \leanok
    Let $x \in \mathbb{R}$ and $n \in \mathbb{N}$. A \textbf{geometric sum}
    is a sum of the form
    \[
        \sum_{k=0}^n x^k.
    \]
\end{definition}

\begin{proposition}
    \label{prop:geom_sum_base_neq_one}
    \lean{GeometricSum.base_ne_one}
    \leanok
    \uses{def:geom_sum}
    Let $x \neq 1$ and $n \in \mathbb{N}$. Then the following inequality 
    holds:
    \[
        \sum_{k=0}^n x^k \leq \frac{x^{n+1} - 1}{x - 1}
    \]
\end{proposition}

\begin{proof}
    \leanok
    We prove the proposition by induction on $n$. The base case is simple:
    \begin{align*}
        \frac{x^{0+1} - 1}{x-1} &= \\
        &= \frac{x^1 - 1}{x-1} \\
        &= \frac{x-1}{x-1} \\
        &= 1 \\
        &= x^0 \\
    \end{align*}
    The inductive step is also straightforward:
    \begin{align*}
        \sum_{k=0}^{n} x^k &= \sum_{k=0}^{n-1} x^k + x^n \\
        &= \frac{x^n - 1}{x-1} + \frac{x-1}{x-1} x^n \\
        &= \frac{x^n - 1}{x-1} + \frac{x^{n+1} - x^n}{x-1} \\
        &= \frac{x^{n+1} - 1}{x-1}
    \end{align*}
\end{proof}

\begin{proposition}
    \label{prop:geom_sum_base_lt_one}
    \lean{GeometricSum.le_of_base_pos_lt_one}
    \leanok
    \uses{def:geom_sum}
    Let $0 < x < 1$ and $n \in \mathbb{N}$. Then the following inequality 
    holds:
    \[
        \sum_{k=0}^n x^k \leq \frac{1}{1-x}
    \]
\end{proposition}

\begin{proof}
    \leanok
    \uses{prop:geom_sum_base_neq_one}
    The inequality follows directly from the previous proposition and the 
    restriction on values of $x$:
    \begin{align*}
        \sum_{k=0}^n x^k &= \frac{x^{n+1} - 1}{x-1} \\
        &= \frac{1 - x^{n+1}}{1-x} \\
        &\leq \frac{1 - x^{n+1}}{1-x}
    \end{align*}
\end{proof}


%%FILE blueprint/src/master_theorem.tex
\chapter{The Master theorem}

Analyzing the time complexity of algorithms, especially recursive ones, is more often 
than not a non-trivial task. For a recursive algorithm, its time complexity can be 
written as a recurrence formula, which is generally not easy, sometimes even impossible 
to solve with a closed formula. In some cases, though, we can find asymptotic bounds 
of the solution, despite not being able to necessarily find the 
precise solution to the recurrence. One large class of such cases is the class of 
divide-and-conquer algorithms, i.e. algorithms that recursively split the
problem into smaller, similarly-sized subproblems. The Master theorem puts asymptotic 
bounds on divide-and-conquer recurrences.

\begin{theorem}[Master theorem for divide-and-conquer recurrences]
    Let $T f : \mathbb{N} \rightarrow \mathbb{N}$ be functions such that the recurrence
    \[
        T(n) = a T(n/b) + f(n)
    \]
    holds for some $a > 0$ and $b > 1$. Let also $f(n) \in \Theta(n^d)$ for some $d \ge 1$. 
    Then the following holds:
    \begin{enumerate}
        \item if $a < b^d$, then $T(n) \in \Theta(n^d)$,
        \item if $a = b^d$, then $T(n) \in \Theta(n^d \log_b{n})$ and
        \item if $a > b^d$, then $T(n) \in \Theta(n^{\log_b{a}})$.
    \end{enumerate}
    In the case where $f$ is bounded above by $n^d$ only above or only below,
    $T$ is also bounded only above or only below by the respective function.
\end{theorem}

We prove this theorem by proving all of its cases separately. 

\begin{definition}
    \label{def:upper_master_rec}
    \lean{UpperMasterRec}
    \leanok
    \uses{def:big_o}
    Let $T f : \mathbb{N} \rightarrow \mathbb{N}$ be functions such that $T$ 
    is monotone and the recurrence
    \[
        T(n) \leq a T(\lceil n/b \rceil) + f(n)
    \]
    holds for all $n \ge n\_0$ for some $a > 0$, $n\_0 > 1$ and $b > n\_0$. 
    Let also $f(n) \in O(n^d)$ for some $d \ge 1$. The above recurrence 
    is an \textbf{upper master recurrence} with parameters $(a, n\_0, b, d)$.
\end{definition}

\begin{definition}
    \label{def:lower_master_rec}
    \lean{LowerMasterRec}
    \leanok
    \uses{def:big_omega}
    Let $T f : \mathbb{N} \rightarrow \mathbb{N}$ be functions such that $T$ 
    is monotone and the recurrence
    \[
        T(n) \geq a T(\lfloor n/b \rfloor) + f(n)
    \]
    holds for all $n \ge n\_0$ for some $a > 0$, $n\_0 > 1$ and $b > n\_0$. 
    Let also $f(n) \in \Omega(n^d)$ for some $d \ge 1$. The above recurrence 
    is a \textbf{lower master recurrence} with parameters $(a, n\_0, b, d)$.
\end{definition}

\begin{lemma}
    \label{lemma:big_o_geom}
    \lean{UpperMasterRec.O_geom}
    \leanok
    \uses{def:upper_master_rec, def:big_o}
    Let $T$ and $f$ form a lower master recurrence with parameters 
    $(a, n\_0, b, d)$. Then $T(n) \in O(n^{\log_b{a}} + 
    \sum_{k=0}^{\lfloor \log_b{n} \rfloor} (\frac{a}{b^d})^k n^d)$.
\end{lemma}

\begin{proof}
    \leanok
    \uses{lemma:big_o_trans, lemma:big_o_add}
    To eliminate ceilings in the recurrence, we substitute $T(n)$ with 
    $S(n) = T(n+b)$. Since $T$ is monotone, $T(n) \leq S(n) = T(n+b)$ holds 
    for all $n$. Therefore, an upper bound on $S$ is also an upper bound on 
    $T$. We must first show that $S$ follows the recurrence of $T$ without 
    ceilings. By the assumption that $b$ is a natural number such that $b > 1$, 
    we have the inequality 
    \begin{align*}
        \lceil \frac{n+b}{b} \rceil &\leq \frac{n+b}{b} + 1 \\
                                    &=    \frac{n}{b} + 2 \\
                                    &\leq \frac{n}{b} + b.
    \end{align*}
    Therefore, by monotonicity of $T$ we have 
    \begin{align*} 
        S(n)  &= T(n+b) \\
              &\leq a T(\frac{n}{b} + b) + f(n+b) \\
              &= a S(\frac{n}{b}) + f(n+b),
    \end{align*}
    which captures the wanted recurrence. Integer division here is defined as 
    the floor of real-number division.
    By substituting the inequality into itself repeatedly, we get
    \begin{align}
        \label{eq:s_rec}
        S(n) &\leq a S(\frac{n}{b}) + f(n + b) \\
             &\leq a S(\frac{n}{b}) + C n^d \\
             &\leq a^2 S(\frac{n}{b^2}) + C n^d + C \frac{a}{b^d} n^d \\
             &\leq a^3 S(\frac{n}{b^3}) + C (1 + \frac{a}{b^d} + 
                                            (\frac{a}{b^d})^2) n^d \\
             &\leq \dots \\
             &\leq a^k S(\frac{n}{b^k}) + C \sum_{i=0}^k 
                                                    (\frac{a}{b^d})^i n^d \\
    \end{align}
    Let $N$ be the integer such that for all $n \geq N \geq n\_0$, the 
    inequality $f(n) \leq C n^d$ holds. Such $N$ exists because 
    $f(n) \in O(n^d)$. We set $k = \lfloor \log_b{\frac{n}{N}} \rfloor$. 
    This choice of $k$ allows the inequality $n \geq N * b^k$ to hold for 
    large enough $n$. 

    Consider both sum parts on the right side of the Equation \ref{eq:s_rec}.
    For the first part, we notice that $n^{\log_b{a}} = a^{\log_b{n}}$.
    Expand $k$ in the exponent and apply monotonicity:
    \begin{align*}
        a^k &= a^{\lfloor \log_b{\frac{n}{N}} \rfloor} \\
            &\leq a^{\log_b{\frac{n}{N}}} \\
            &\leq a^{\log_b{n}}
    \end{align*}
    This implies $a^k \in O(n^{\log_b{a}})$. We also need $S(\frac{n}{b^k})$ to
    be bounded by a constant. We show that $S(\frac{n}{b^k}) \leq S(N*b)$.
    Since $S$ is monotone, this is equivalent to showing 
    $\frac{n}{b^k} \leq N*b$. We rewrite the left side as
    \begin{align*}
        \frac{n}{b^k} &= \frac{n}{b^{\lfloor \log_b{\frac{n}{N}} \rfloor}} \\
                      &= \frac{n}{b^{\lfloor \log_b{n} - \log_b{N} \rfloor}} \\
                      &= \frac{b^{\log_b{n}}}
                              {b^{\lfloor \log_b{n} - \log_b{N} \rfloor}} \\
                      &= b^{\log_b{n} - 
                            \lfloor \log_b{n} - \log_b{N} \rfloor}
    \end{align*}
    Then, rewrite the right side as $N*b = b^{\log_b{N} + 1}$. With both 
    sides written as exponents of $b$, we only need to prove
    \[\log_b{n} - \lfloor \log_b{n} - \log_b{N} \rfloor \leq \log_b{N} + 1.\]
    By swapping terms, we get
    \[\log_b{n} - \log_b{N} \leq \lfloor \log_b{n} - \log_b{N} \rfloor + 1,\]
    which holds for all real numbers. 

    For the second part, the inequality $k \leq \lfloor \log_b{n} \rfloor$ 
    holds by monotonicity of logarithms. As geometric sums are monotone in 
    the exponent, we get $\sum_{i=0}^k (\frac{a}{b^d})^i n^d \in 
    O(\sum_{i=0}^{\lfloor \log_b{n} \rfloor} (\frac{a}{b^d})^i n^d)$.
\end{proof}

\begin{theorem}
    \label{thm:big_o_of_lt}
    \lean{UpperMasterRec.O_of_lt}
    \leanok
    \uses{def:upper_master_rec, def:big_o}
    Let $T$ and $f$ form a lower master recurrence with parameters 
    $(a, n\_0, b, d)$, where $a < b^d$. Then $T(n) \in O(n^d)$.
\end{theorem}

\begin{proof}
    \leanok
    \uses{lemma:big_o_geom, lemma:big_o_trans, lemma:big_o_add}
    First, we apply Lemma \ref{lemma:big_o_geom}. Since $a < b^d$, we have 
    $\frac{a}{b^d} < 1$. By basic properties of geometric sums, we get
    \begin{align*}
        T(n) &\leq \sum_{i=0}^{\lfloor \log_b{n} \rfloor} 
                        (\frac{a}{b^d})^i n^d \\
             &\leq \frac{1}{1 - \frac{a}{b^d}} n^d,
    \end{align*}
    which proves the upper bound.
\end{proof}

\begin{theorem}
    \label{thm:big_o_of_eq}
    \lean{UpperMasterRec.O_of_eq}
    \leanok
    \uses{def:upper_master_rec, def:big_o}
    Let $T$ and $f$ form an upper master recurrence with parameters
    $(a, n\_0, b, d)$, where $a = b^d$. Then $T(n) \in O(n^d \log_b{n})$.
\end{theorem}

\begin{proof}
    \leanok
    \uses{lemma:big_o_geom, lemma:big_o_trans, lemma:big_o_add}
    After applying Lemma \ref{lemma:big_o_geom}, we note that $\log_b{a} = d$
    and then the proof boils down to showing that the geometric sum is bounded 
    by $\log_b{n}$. Since $\frac{a}{b^d} = 1$, the geometric sum equals
    $\lfloor \log_b{n} \rfloor$, which is obviously bounded by $\log_b{n}$.
\end{proof}

\begin{theorem}
    \label{thm:big_o_of_gt}
    \lean{UpperMasterRec.O_of_gt}
    \leanok
    \uses{def:upper_master_rec, def:big_o}
    Let $T$ and $f$ form an upper master recurrence with parameters
    $(a, n\_0, b, d)$, where $a > b^d$. Then $T(n) \in O(n^{\log_b{a}})$.
\end{theorem}

\begin{proof}
    \leanok
    \uses{lemma:big_o_geom, lemma:big_o_trans, lemma:big_o_add}
    After applying Lemma \ref{lemma:big_o_geom}, the left side of the sum if 
    trivially bounded by $n^{\log_b{a}}$. We are left with the right summand, 
    which we transform using an inequality involving the geometric sum:
    \begin{align*}
        \sum_{k=0}^{\lfloor \log_b{n} \rfloor} (\frac{a}{b^d})^k n^d
            &\leq \frac{1}{\frac{a}{b^d} - 1}
                    ((\frac{a}{b^d})^{\lfloor \log_b{n} \rfloor} - 1) n^d \\
            &\leq \frac{1}{\frac{a}{b^d} - 1}
                    ((\frac{a}{b^d})^{\log_b{n}} - 1) n^d \\
            &= \frac{1}{\frac{a}{b^d} - 1}
                    (b^{\log_b{\frac{a}{b^d}}})^{\log_b{n}} n^d \\
            &= \frac{1}{\frac{a}{b^d} - 1}
                    (b^{\log_b{n}})^{\log_b{\frac{a}{b^d}}} n^d \\
            &= \frac{1}{\frac{a}{b^d} - 1}
                    n^{\log_b{a} - d} n^d \\
            &= \frac{1}{\frac{a}{b^d} - 1}
                    \frac{n^{\log_b{a}}}{n^d} n^d \\
            &= \frac{1}{\frac{a}{b^d} - 1}
                    n^{\log_b{a}} \\
    \end{align*}
    Since $a > b^d$, $\frac{a}{b^d} > 1$ holds, so 
    $\frac{1}{\frac{a}{b^d} - 1} > 0$, which proves boundedness by 
    $n^{\log_b{a}}$.
\end{proof}


\begin{lemma}
    \label{lemma:big_omega_geom}
    \lean{LowerMasterRec.Ω_geom}
    \leanok
    \uses{def:lower_master_rec, def:big_omega}
    Let $T$ and $f$ form an upper master recurrence with parameters 
    $(a, n\_0, b, d)$. Then $T(n) \in 
    \Omega(\sum_{k=0}^{\lfloor \log_b{n} \rfloor} (\frac{a}{b^d})^k n^d)$.
\end{lemma}

\begin{proof}
    \leanok
    \uses{lemma:big_omega_trans}
    We consider the recurrence formula with ceilings replaced by floors. If the
    resulting inequality holds, then so does the master recurrence, so it 
    suffices to prove the lower bound for this inequality.

    By substituting the inequality into itself repeatedly, we get
    \begin{align}
        \label{eq:t_rec}
        T(n) &\leq a T(\frac{n}{b}) + f(n + b) \\
             &\leq a T(\frac{n}{b}) + C n^d \\
             &\leq a^2 T(\frac{n}{b^2}) + C n^d + C \frac{a}{b^d} n^d \\
             &\leq a^3 T(\frac{n}{b^3}) + C (1 + \frac{a}{b^d} + 
                                            (\frac{a}{b^d})^2) n^d \\
             &\leq \dots \\
             &\leq a^k T(\frac{n}{b^k}) + C \sum_{i=0}^k 
                                                    (\frac{a}{b^d})^i n^d \\
    \end{align}
    We set $k = \lfloor \log_b{n} \rfloor$. This choice of $k$ allows the 
    inequality $n \geq b^k$ to hold for large enough $n$. Here $C$ is a 
    positive constant such that $f(n) \geq C n^d$ for all $n \geq n\_0$. 
    Such a constant exists, because $f(n) \in \Omega(n^d)$ implies
    $f(n) \geq C\_0 n^d$ for some $C\_0 > 0$ for all $n \geq N$ for some $N$.
    For $n\_0 \leq n \leq N$, The argument is as follows. The set of natural 
    numbers between $n\_0$ and $N$ is finite, so the image of $f$ on this
    set has a maximal element $M$. We then have $f(n) \geq \frac{M}{N^d} n^d$.

    The second summand in the right side of Equation \ref{eq:t_rec} is
    trivially bounded above by 
    $\sum_{k=0}^{\lfloor \log_b{n} \rfloor} (\frac{a}{b^d})^k n^d$. This is 
    sufficient to prove the upper bound od $T(n)$ as the left summand is 
    non-negative.
\end{proof}

\begin{lemma}
    \label{lemma:big_omega_pow}
    \lean{LowerMasterRec.Ω_pow_d}
    \leanok
    \uses{def:lower_master_rec, def:big_omega}
    Let $T$ and $f$ form a lower master recurrence with parameters 
    $(a, n\_0, b, d)$. Then $T(n) \in \Omega(n^d)$.
\end{lemma}

\begin{proof}
    \leanok
    \uses{lemma:big_omega_trans}
    Since $T(n) \geq f(n)$ for all $n \ge n\_0$, the lower bound follows 
    directly from $f(n) \in \Omega(n^d)$.
\end{proof}

\begin{theorem}
    \label{thm:big_omega_of_eq}
    \lean{LowerMasterRec.Ω_of_eq}
    \leanok
    \uses{def:lower_master_rec, def:big_omega}
    Let $T$ and $f$ form a lower master recurrence with parameters 
    $(a, n\_0, b, d)$ where $a = b^d$. Then $T(n) \in \Omega(n^d \log_b{n})$.
\end{theorem}

\begin{proof}
    \leanok
    \uses{lemma:big_omega_geom, lemma:big_omega_trans}
    By Lemma \ref{lemma:big_omega_geom}, it suffices to show that
    $\sum_{i=0}^{\lfloor \log_b{n} \rfloor} (\frac{a}{b^d})^i \in 
    \Omega(\log_b{n})$. Applying equality $a = b^d$, the sum simplifies to
    \begin{align*}
        \sum_{i=0}^{\lfloor \log_b{n} \rfloor} (\frac{a}{b^d})^i &= 
            \sum_{i=0}^{\lfloor \log_b{n} \rfloor} 1^i &=
            \lfloor \log_b{n} \rfloor.
    \end{align*}
\end{proof}

\begin{theorem}
    \label{thm:big_omega_of_gt}
    \lean{LowerMasterRec.Ω_of_gt}
    \leanok
    \uses{def:lower_master_rec, def:big_omega}
    Let $T$ and $f$ form an upper master recurrence with parameters 
    $(a, n\_0, b, d)$ where $a > b^d$. Then $T(n) \in \Omega(n^{\log_b{a}})$.
\end{theorem}

\begin{proof}
    \leanok
    \uses{lemma:big_omega_geom, lemma:big_omega_trans}
    By Lemma \ref{lemma:big_omega_geom}, we need to show that 
    $\sum_{i=0}^{\lfloor \log_b{n} \rfloor} (\frac{a}{b^d})^i n^d
    \in \Omega(n^{\log_b{a}})$. By $a > b^d$, we have 
    \begin{align*}
        \sum_{i=0}^{\lfloor \log_b{n} \rfloor} (\frac{a}{b^d})^i n^d 
        &\geq (a b^{-d} - 1)^{-1}
            ((a b^{-d})^{\lfloor \log_b{n} \rfloor} - 1) n^d \\
        &\geq 2^{-1} (a b^{-d} - 1)^{-1}
            (a b^{-d})^{\lfloor \log_b{n} \rfloor} n^d \\
        &\geq 2^{-1} (a b^{-d} - 1)^{-1} (a b^{-d})^{\log_b{n} - 1} n^d \\
        &\geq 2^{-1} a^{-1} b^d (a b^{-d} - 1)^{-1} a^{\log_b{n}} 
            (b^{\log_b{n} - 1})^{-d} n^d \\
        &\geq 2^{-1} a^{-1} b^d (a b^{-d} - 1)^{-1} n^{\log_b{a}} n^{-d} n^d \\
        &\geq 2^{-1} a^{-1} b^d (a b^{-d} - 1)^{-1} n^{\log_b{a}},
    \end{align*}
    which proves the bound.
\end{proof}


\begin{corollary}
    \label{thm:theta_of_lt}
    \lean{MasterRec.Θ_of_lt}
    \leanok
    \uses{def:upper_master_rec, def:lower_master_rec, def:theta}
    Let $T$ and $f$ form an upper and lower master recurrence with parameters 
    $(a, n\_0, b, d)$ where $a < b^d$. Then $T(n) \in \Theta(n^d)$.
\end{corollary}

\begin{proof}
    \leanok
    \uses{thm:big_o_and_omega_iff_theta, thm:big_o_of_lt, lemma:big_omega_pow}
    By Theorem \ref{thm:big_o_and_omega_iff_theta}, it suffices to show
    lower and upper bounds for $T$, which we already proved in Theorem
    \ref{thm:big_o_of_lt} and Lemma \ref{lemma:big_omega_pow}.
\end{proof}

\begin{corollary}
    \label{thm:theta_of_eq}
    \lean{MasterRec.Θ_of_eq}
    \leanok
    \uses{def:upper_master_rec, def:lower_master_rec, def:theta}
    Let $T$ and $f$ form an upper and lower master recurrence with parameters 
    $(a, n\_0, b, d)$ where $a = b^d$. Then $T(n) \in \Theta(n^d \log_b{a})$.
\end{corollary}

\begin{proof}
    \leanok
    \uses{thm:big_o_and_omega_iff_theta, thm:big_o_of_eq, thm:big_omega_of_eq}
    By Theorem \ref{thm:big_o_and_omega_iff_theta}, it suffices to show
    lower and upper bounds for $T$, which we already proved in Theorems
    \ref{thm:big_o_of_eq} and \ref{thm:big_omega_of_eq}.
\end{proof}

\begin{corollary}
    \label{thm:theta_of_gt}
    \lean{MasterRec.Θ_of_gt}
    \leanok
    \uses{def:upper_master_rec, def:lower_master_rec, def:theta}
    Let $T$ and $f$ form an upper and lower master recurrence with parameters 
    $(a, n\_0, b, d)$ where $a > b^d$. Then $T(n) \in \Theta(n^{\log_b{n}})$.
\end{corollary}

\begin{proof}
    \leanok
    \uses{thm:big_o_and_omega_iff_theta, thm:big_o_of_gt, thm:big_omega_of_gt}
    By Theorem \ref{thm:big_o_and_omega_iff_theta}, it suffices to show
    lower and upper bounds for $T$, which we already proved in Theorems
    \ref{thm:big_o_of_gt} and \ref{thm:big_omega_of_gt}.
\end{proof}

